\documentclass[letterpaper]{article}

\usepackage[letterpaper,margin=1in]{geometry}

% Compile from the repository root so that this path resolves to
% the pinned WSP submodule.
\input{wsp/testing/test-case-library.tex}

\begin{document}

\def\TCID{%
  TC-0009
}

\def\TCTitle{%
  Library DLL: exports, version and registration
}

\def\TCPurpose{%
  Verify the shared library \texttt{wfcmscomctl.dll} as any program uses it:
  the COM entry points, the version resource, the registration round trip, and
  the controls it creates through COM (ranges, collections, events, user
  input).
}

\def\TCPriority{%
  \TCPriorityHigh
}

\def\TCRequirementRef{%
  REQ-0025
}

\def\TCDesignRef{%
  Not applicable; the rules are in the named header.
}

\def\TCFeatureRef{%
  \texttt{tests/mscomctl\_dll\_test.c}, \texttt{tests/dll\_check.h},
  \texttt{src/registration.cpp}, \texttt{src/com\_behaviour.hpp},
  \texttt{src/com\_collections.hpp}, \texttt{src/com\_input.hpp} and the
  runner \texttt{tests/tc-0009-dll.ps1}
}

\def\TCTestTechnique{%
  Black-box integration testing through the COM interfaces (ISO/IEC/IEEE
  29119-4 specification-based techniques).
}

\def\TCPreconditions{%
  Windows; the library built with \texttt{WFC\_MSCOMCTL\_SHARED=ON} so that
  \texttt{wfcmscomctl.dll} and the C test \texttt{mscomctl\_dll\_test} exist.
  The current user's registry may be written (\texttt{HKCU}).
}

\def\TCEnvironment{%
  Windows x86, x64 or ARM64; CTest runs the test as \texttt{TC-0009-dll}
  (or \texttt{wfc-mscomctl-TC-0009-dll} when WFC adds the library) with
  resource lock \texttt{registry} and treats exit code 77 as skipped.
}

\def\TCAssumptions{%
  Neither the library nor the Microsoft class ids are already registered;
  otherwise the test leaves the registry alone and exits with 77.
}

\def\TCInputData{%
  The path of the DLL, its version (\texttt{VERSION}) and its target
  (\texttt{TARGET}) as arguments; the class ids of the nine controls.
}

\def\TCInitialState{%
  No registration of the library exists in the registry of the current user.
}

\def\TCProcedure{%
  \begin{enumerate}
    \item Load the DLL and look for the four COM entry points; compare the
      version resource with the version and target given.
    \item Register the library under WFC's ProgIDs and class ids, check every
      class, unregister it, and repeat the round trip under the Microsoft
      ProgIDs.
    \item Create a Slider and a ProgressBar through COM and check the
      defaults, error 380 for a Value outside the range and the Value brought
      inside a narrower range.
    \item Create TabStrip, StatusBar and Toolbar, ListView, ImageList,
      TreeView and ImageCombo objects and exercise their collections (Add,
      Count, Item, Remove, For Each), the library errors, the events and the
      input of the user (\texttt{WfcInput}).
    \item Create the TabStrip once more and call a property and a method.
  \end{enumerate}
}

\def\TCExpectedResult{%
  \begin{enumerate}
    \item Every entry point exists and the version resource matches the
      version and target.
    \item Registration and unregistration succeed for both sets of names and
      leave the registry as it was.
    \item The controls answer as the rules of the library say: the numbered
      errors (380, 35600, 35602, 35609) appear where expected, the first tab
      is selected, collections keep order and the events and input arrive in
      order.
    \item The program exits with code 0 (or 77 when it skipped).
  \end{enumerate}
}

\def\TCPostConditions{%
  The registry of the current user is as before the test.
}

\TCTemplate

\end{document}
